\documentclass{article}
\usepackage[utf8]{inputenc}
\usepackage[absolute]{textpos}
\usepackage[default]{raleway}
\usepackage{titlesec, comment, tabularx, makecell, listings, array, setspace, geometry, graphicx, xcolor, xparse, fancyvrb, relsize, fancyhdr, booktabs, multirow, hyperref}
\usepackage{colortbl}
%\geometry{a4paper, left=2cm, right=2cm, top=2cm, bottom=2.5cm}
\renewcommand{\headrulewidth}{0pt}

% Definisci uno stile per i comandi git
\definecolor{light-gray}{gray}{0.92}

\lstdefinestyle{code}{
    frame=single,
    framesep=1mm,
    rulecolor=\color{light-gray},
    backgroundcolor=\color{light-gray},
    basicstyle=\ttfamily,
}

% ----------------------------- Definizione tabella ---------------------------

\newcolumntype{C}[1]{>{\centering\arraybackslash}m{#1}}

%\setcellgapes{2ex} % Imposta l'altezza dell'header (2ex)


% ------------------------------Metadati indice --------------------------------
\title{\textbf{\fontsize{28}{6}\selectfont Indice}}
\author{\fontsize{14}{6}\selectfont ByteOps}
\date{Marzo 01, 2024}


% -----------------------------Creazione footer --------------------------------

\pagestyle{fancy}
\fancyhf{}
\renewcommand{\footrulewidth}{0.4pt}
\lfoot{
    \parbox[c]{2cm}{\includegraphics[width=2cm]{../../Images/logo.png}}
    \textcolor[RGB]{120, 120, 120}{$\cdot$ Verbale Interno}
}
\rfoot{\thepage}

% --------------------------Modifica formato hyperlinks ------------------------

\hypersetup{
    colorlinks=true,
    linkcolor=black,
    filecolor=black,      
    pdftitle={Verbale Interno 01/03/2024},  %inserisci data verbale
    pdfpagemode=FullScreen,
}

% ------------------------------- Valore sotto-paragrafi indice --------------------------------------

\setcounter{secnumdepth}{4}
\setcounter{tocdepth}{4}

\titleformat{\section}
{\normalfont\huge\bfseries}{\thesection}{0.2cm}{}
\titlespacing*{\paragraph}{0pt}{0.5cm}{0.1cm}

\titleformat{\subsection}
{\normalfont\Large\bfseries}{\thesubsection}{0.2cm}{}
\titlespacing*{\paragraph}{0pt}{0.5cm}{0.1cm}

\titleformat{\subsubsection}
{\normalfont\large\bfseries}{\thesubsubsection}{0.2cm}{}
\titlespacing*{\paragraph}{0pt}{0.5cm}{0.1cm}

\titleformat{\paragraph}
{\normalfont\normalsize\bfseries}{\theparagraph}{0.2cm}{}
\titlespacing*{\paragraph}{0pt}{0.5cm}{0.1cm}

% ------------------------------- Front Page ---------------------------------------

\begin{document}

% --------------------------Aggiunta firma finale ------------------------
\begin{textblock*}{\textwidth}(0.85\textwidth, 1.16\textheight)
    Il responsabile: R. Smanio
\end{textblock*}
% ------------------------------------------------------------------------

\pagestyle{fancy}
\begin{center}
\includegraphics[width = 0.7\textwidth]{../../Images/logo.png} \\
\vspace{0.2cm}
\textcolor[RGB]{60, 60, 60}{\textit{ByteOps.swe@gmail.com}} \\
\vspace{1cm}
\fontsize{16}{6}\selectfont Verbale Interno $\cdot$ Data: 01/03/2024 \\
\vspace{0.5cm}
\end{center}

\section*{Informazioni documento}
\def\arraystretch{1.2}
\begin{tabular}{>{\raggedleft\arraybackslash}p{0.3\textwidth}|>{\raggedright\arraybackslash}p{0.6\textwidth}c}
\hline
\addlinespace
\textbf{Luogo} & Discord \vspace{10pt} \\
\textbf{Orario} & 17:30 - 18:15 \vspace{10pt} \\
\textbf{Redattore} & N. Preto \vspace{10pt} \\
\textbf{Verificatore} & A. Barutta \vspace{10pt} \\
\textbf{Amministratore} & E. Hysa \vspace{10pt} \\
\textbf{Destinatari} & T. Vardanega \\ & R. Cardin \vspace{10pt} \\
\multirow[t]{7}{*}{\textbf{Partecipanti}} & A. Barutta \\ & E. Hysa \\ & R. Smanio \\ & D. Diotto \\ & F. Pozza \\ & L. Skenderi \\ & N. Preto \vspace{10pt} \\
\end{tabular}
\pagebreak 

% ------------------------- Changelog ----------------------------

\section*{Registro delle modifiche}

\begin{tabular}{|C{2.5cm}|C{2.5cm}|C{2.5cm}|C{2.5cm}|C{2.5cm}|}
    \hline
    \textbf{Versione} & \textbf{Data} & \textbf{Autore} & \textbf{Verificatore} & \textbf{Dettaglio} \\
    \hline \hline
    0.0.2 & 02/03/2024 & N. Preto & A. Barutta & Aggiunta issue mancanti \\
    \hline
    0.0.1 & 01/03/2024 & N. Preto & A. Barutta & Redazione documento \\
    \hline
\end{tabular}
\pagebreak

% ------------------------- Generazione automatica indice ----------------------
\setstretch{1.5}
\maketitle
\thispagestyle{fancy}
\tableofcontents
\setstretch{1.2}
\pagebreak

% ------------------------ INIZIO DOCUMENTO ----------------------
\flushleft

\section{Revisione del periodo precedente}
Nella settimana successiva all'ultimo meeting interno, è stato integrato con successo Faust nel prodotto \textit{software}\textsubscript{\textit{G}}, rispondendo efficacemente alle criticità evidenziate dal Professor Cardin durante la revisione RTB e garantendo l'aderenza ai principi di ingegneria del software. \\
Successivamente all'\textit{integrazione}\textsubscript{\textit{G}} di Faust, le principali attività a cui si sono dedicati i membri del team riguardano lo sviluppo dei test e la redazione del \textit{Manuale Utente}.
Inoltre, l'Amministratore si è dedicato al miglioramento della struttura e dell'organizzazione del \textit{repository}\textsubscript{\textit{G}} GitHub, seguendo i suggerimenti ricevuti dal Professor Vardanega nella valutazione della revisione RTB. \\
In aggiunta, sono state apportate modifiche all'aspetto di alcuni pannelli della \textit{dashboard}\textsubscript{\textit{G}} \textit{Grafana}\textsubscript{\textit{G}}, aumentando così la qualità dell'interfaccia e dell'user experience.

\section{Ordine del giorno}
    \subsection{Discussione sulle difficoltà emerse nell'attività di testing}
    All'inizio dell'incontro, è stata posta l'attenzione sui \textit{test}\textsubscript{\textit{G}} che sono attualmente in fase di sviluppo. \\
    Nei giorni precedenti infatti, sono state riscontrate difficoltà nella realizzazione dei \textit{test}\textsubscript{\textit{G}} di integrazione tra il producer e il server Kafka, principalmente a causa della poca esperienza del team in tale ambito. \\
    A seguito delle discussioni tenute durante il SAL odierno, l'azienda proponente ha proposto di fornirci supporto per la risoluzione dei problemi riscontrati. Inoltre, si è deciso di fare pratica con dei minimal working examples in modo tale da acquisire le competenze e l'esperienza necessaria per poter realizzare i test di integrazione nel modo corretto in un contesto più complesso come quello del progetto.

    \subsection{Discussione sul livello di dettaglio del Manuale Utente}
    Dopo un'attenta riflessione sul livello di dettaglio da mantenere riguardo alle istruzioni e alle procedure nel \textit{Manuale Utente}, si era inizialmente optato per una descrizione più semplice delle funzionalità offerte e delle relative modalità di utilizzo. Questa scelta mirava a rendere la lettura e la comprensione del manuale accessibili, evitando complicazioni. Tuttavia, una successiva riflessione più approfondita del team ha portato alla conclusione che fosse necessario ampliare ulteriormente le descrizioni delle funzionalità e fornire indicazioni più dettagliate sul loro scopo e sul modo migliore per utilizzarle. In quest'ottica, diventa fondamentale trovare un equilibrio tra precisione, dettaglio e approfondimento, mantenendo allo stesso tempo la semplicità di lettura e comprensione. \\
    L'adozione di una maggior granularità delle informazioni mira a garantire una comprensione completa degli aspetti e dei dettagli meno evidenti del prodotto, riducendo al contempo il rischio di fraintendimenti e ambiguità. L'obiettivo finale è rendere il Manuale Utente più accessibile, preciso e utile per il lettore, fornendo informazioni dettagliate sulle modalità di utilizzo di tutte le funzionalità offerte, soprattutto su quelle che potrebbero non risultare immediate o intuitive.

    \subsection{Discussione riguardante il sistema di autenticazione}
    Si è evidenziato un ulteriore punto di interesse riguardante il mantenimento del \textit{sistema}\textsubscript{\textit{G}} di autenticazione fornito da \textit{Grafana}\textsubscript{\textit{G}} e, dopo aver analizzato attentamente le diverse opinioni espresse in merito, si è concluso che mantenere l'autenticazione è fondamentale poiché garantisce un livello di sicurezza superiore. \\
    È da sottolineare che, come concordato con la proponente, non saranno previsti account con differenti livelli di privilegio all'interno del software. Ciò è dovuto al fatto che gli utenti finali, ovvero i membri delle autorità locali, devono avere accesso completo a tutte le funzionalità del sistema. \\
    In conclusione, per accedere alla dashboard, gli utenti facenti parte delle autorità locali non dovranno registrarsi in autonomia, ma verranno fornite loro delle credenziali di accesso (nome utente e password) tramite un canale affidabile.

    \subsection{Controllo conformità ai requisiti}
    È stato riconosciuto l'importante bisogno di condurre un'attenta revisione dei requisiti delineati nel documento "\textit{Analisi dei Requisiti}". Ciò si configura come un passaggio cruciale per garantire che le modifiche apportate alla dashboard siano conformi ai requisiti stabiliti all'inizio del progetto. L'obiettivo principale è garantire che il gruppo non abbia introdotto elementi superflui o deviazioni dalle specifiche inizialmente definite.

    \section{Attività da svolgere}
    \begin{center}
        \begin{tabular}{|C{7cm}|C{1,5cm}|C{3cm}|}
            \hline
            \textbf{Titolo} & \textbf{\# Issue} & \textbf{Verificatore} \\
            \hline
            \hline
            Approfondimento paragrafi Manuale Utente & 96 & A. Barutta \\
            \hline
            Ricontrollo AdR dopo prima stesura MU & 97 & A. Barutta \\
            \hline
            Redazione verbale interno 01/03/2024 & 99 & N. Preto \\
            \hline
            Stesura periodo 10 - PdP & 100 & A. Barutta \\
            \hline
            Sviluppo \textit{test}\textsubscript{\textit{G}} di unità sensori & 101 & R. Smanio \\
            \hline
            Sviluppo \textit{test}\textsubscript{\textit{G}} di \textit{integrazione}\textsubscript{\textit{G}} \textit{kafka}\textsubscript{\textit{G}} - \textit{python}\textsubscript{\textit{G}} & 102 & R. Smanio \\
            \hline
            Sviluppo \textit{test}\textsubscript{\textit{G}} di \textit{integrazione}\textsubscript{\textit{G}} \textit{Kafka}\textsubscript{\textit{G}} - \textit{Python}\textsubscript{\textit{G}} - \textit{Clickhouse}\textsubscript{\textit{G}} & 103 & A. Barutta \\
            \hline
        \end{tabular}
    \end{center}
\end{document}
